% =============================================================================
% GATEWAY ARCHITECTURE
% =============================================================================

\newpage
\section{Gateway Architecture}
\label{sec:gateway-architecture}

This section documents the \texttt{star-gateway} Go service that bridges the Raspberry Pi 5 control stack with the ESP32-S3 motor controller.

\subsection{Overview}

The gateway service runs on the Raspberry Pi 5 and serves as the communication bridge between:

\begin{itemize}[leftmargin=1.5em,topsep=4pt,itemsep=2pt]
    \item \textbf{High-level side:} gRPC clients (Nav2/ROS2 for autonomous navigation, UI for manual control)
    \item \textbf{Low-level side:} ESP32-S3 motor controller via SPI at 10 MHz
\end{itemize}

\noindent\makebox[\textwidth]{%
\begin{tikzpicture}[
    font=\small,
    node distance=0.5cm,
    box/.style={
        draw,
        thick,
        rectangle,
        minimum width=3.5cm,
        minimum height=1cm,
        align=center
    }
]
    % High-level clients
    \node[box, fill=blue!20] (nav2) at (0,0) {Nav2/ROS2\\(Autonomous)};
    \node[box, fill=green!20, right=of nav2] (ui) {star-ui\\(Manual Control)};

    % Gateway
    \node[box, fill=orange!30, below=1cm of $(nav2)!0.5!(ui)$, minimum width=8cm] (gateway) {\textbf{star-gateway}\\gRPC Server + SPI Controller};

    % ESP32
    \node[box, fill=cyan!20, below=1cm of gateway, minimum width=8cm] (esp32) {\textbf{ESP32-S3}\\Motor Controller + Sensors};

    % Arrows
    \draw[->, thick] (nav2.south) -- ++(0,-0.3) -| (gateway.north west);
    \draw[->, thick] (ui.south) -- ++(0,-0.3) -| (gateway.north east);
    \draw[<->, thick] (gateway.south) -- node[right] {SPI @ 10 MHz} (esp32.north);

    % Labels
    \node[above=0.1cm of nav2, font=\footnotesize\bfseries] {gRPC Clients};
\end{tikzpicture}%
}

\subsection{Protocol Stack Implementation}

The gateway implements Layers 1--4 of the communication protocol (see Section~\ref{sec:protocol-stack}).

\begin{table}[H]
\centering
\caption{Gateway Protocol Stack Implementation}
\label{tab:gateway-stack}
\begin{tabular}{|c|l|l|l|}
\hline
\textbf{Layer} & \textbf{Name} & \textbf{Go Package} & \textbf{Status} \\
\hline
5 & Application & \texttt{internal/service/} & Placeholder \\
4 & Serialization & \texttt{star-proto/gen/go/} & Generated \\
3 & ARQ & \texttt{internal/arq/} & Placeholder \\
2 & Framing & \texttt{internal/frame/} & Constants defined \\
1 & Transport & \texttt{internal/transport/} & Placeholder \\
\hline
\end{tabular}
\end{table}

\subsection{Directory Structure}

\begin{verbatim}
star-gateway/
├── cmd/
│   └── star-gateway/
│       └── main.go              # Entry point
├── internal/
│   ├── transport/               # Layer 1: SPI transport
│   │   ├── spi.go               # Transport interface
│   │   └── spi_test.go
│   ├── frame/                   # Layer 2: Frame protocol
│   │   ├── frame.go             # Constants and types
│   │   ├── encoder.go           # Frame encoder
│   │   ├── decoder.go           # Frame decoder
│   │   └── frame_test.go
│   ├── arq/                     # Layer 3: ARQ protocol
│   │   ├── arq.go               # Stop-and-Wait ARQ
│   │   └── arq_test.go
│   └── service/                 # Layer 5: gRPC services
│       ├── motor_control.go
│       ├── telemetry.go
│       ├── battery.go
│       ├── configuration.go
│       └── firmware.go
├── go.mod
└── CLAUDE.md
\end{verbatim}

\subsection{Frame Package}

The \texttt{internal/frame} package defines the wire protocol constants and types.

\subsubsection{Frame Constants}

\begin{table}[H]
\centering
\caption{Frame Protocol Constants}
\label{tab:frame-constants}
\begin{tabular}{|l|c|l|}
\hline
\textbf{Constant} & \textbf{Value} & \textbf{Description} \\
\hline
\texttt{SyncWord} & \texttt{0x55AA} & Frame sync marker (alternating bits: \texttt{[0x55, 0xAA]}) \\
\texttt{SyncSize} & 2 bytes & Sync word size \\
\texttt{HeaderSize} & 6 bytes & SEQ(2) + LEN(2) + TYPE(1) + FLAGS(1) \\
\texttt{CRCSize} & 4 bytes & CRC-32 checksum \\
\texttt{MaxPayloadSize} & 1024 bytes & Maximum protobuf payload \\
\texttt{MaxFrameSize} & 1036 bytes & Total frame size \\
\texttt{MinFrameSize} & 12 bytes & Empty payload frame \\
\hline
\end{tabular}
\end{table}

\subsubsection{Frame Types}

\begin{itemize}[leftmargin=1.5em,topsep=4pt,itemsep=2pt]
    \item \texttt{FrameTypeCommand} -- Command from controller to peripheral
    \item \texttt{FrameTypeResponse} -- Response from peripheral to controller
    \item \texttt{FrameTypeAck} -- Acknowledgment of successful receipt
    \item \texttt{FrameTypeNack} -- Negative acknowledgment (CRC error, etc.)
\end{itemize}

\subsection{Transport Package}

The \texttt{internal/transport} package provides the SPI transport layer.

\subsubsection{SPI Configuration}

\begin{table}[H]
\centering
\caption{SPI Configuration Constants}
\label{tab:spi-config}
\begin{tabular}{|l|l|l|}
\hline
\textbf{Parameter} & \textbf{Value} & \textbf{Description} \\
\hline
Device & \texttt{/dev/spidev0.0} & SPI device path \\
Speed & 10 MHz & SPI clock frequency \\
Mode & 0 & CPOL=0, CPHA=0 \\
Bits per word & 8 & Word size \\
Timeout & 100 ms & Default operation timeout \\
\hline
\end{tabular}
\end{table}

\subsubsection{Transport Interface}

\begin{verbatim}
type Transport interface {
    Send(data []byte) (int, error)
    Receive(maxLen int) ([]byte, error)
    Transfer(txData []byte) ([]byte, error)
    Close() error
}
\end{verbatim}

\subsection{ARQ Package}

The \texttt{internal/arq} package implements Stop-and-Wait ARQ for reliable delivery.

\subsubsection{ARQ Parameters}

\begin{table}[H]
\centering
\caption{ARQ Protocol Parameters}
\label{tab:arq-params}
\begin{tabular}{|l|l|l|}
\hline
\textbf{Parameter} & \textbf{Value} & \textbf{Description} \\
\hline
Max Retries & 3 & Maximum transmission attempts \\
Timeout & 10 ms & ACK wait timeout \\
Sequence Max & 0xFFFF & 16-bit wraparound \\
\hline
\end{tabular}
\end{table}

\subsubsection{ARQ State Machine}

\begin{itemize}[leftmargin=1.5em,topsep=4pt,itemsep=2pt]
    \item \textbf{IDLE} -- Ready to send or receive
    \item \textbf{WAITING\_ACK} -- Waiting for acknowledgment
    \item \textbf{RETRANSMITTING} -- Retransmitting after timeout/NACK
    \item \textbf{ERROR} -- Max retries exceeded
\end{itemize}

\subsection{gRPC Services}

The gateway exposes five gRPC services (see Section~\ref{sec:service-specifications} for detailed specifications):

\begin{table}[H]
\centering
\caption{Gateway gRPC Services}
\label{tab:gateway-services}
\begin{tabular}{|l|l|}
\hline
\textbf{Service} & \textbf{Purpose} \\
\hline
MotorControlService & Differential drive control, encoder streaming \\
TelemetryService & IMU, GPS, system status \\
BatteryManagementService & BQ7850 BMS monitoring \\
ConfigurationService & Runtime parameters, NVS persistence \\
FirmwareUpdateService & OTA updates, rollback \\
\hline
\end{tabular}
\end{table}

\subsection{Build and Test}

\begin{verbatim}
# Build the gateway binary
cd star-gateway
go build ./cmd/star-gateway

# Run all tests
go test ./...

# Run with verbose output
go test -v ./...

# Format code
go fmt ./...

# Static analysis
go vet ./...
\end{verbatim}

\subsection{Dependencies}

The gateway depends on:

\begin{itemize}[leftmargin=1.5em,topsep=4pt,itemsep=2pt]
    \item \texttt{github.com/Locked-Inc/star-proto/gen/go} -- Generated protobuf and gRPC code
    \item \texttt{google.golang.org/grpc} -- gRPC framework
    \item \texttt{google.golang.org/protobuf} -- Protocol Buffers runtime
\end{itemize}

The \texttt{go.work} file at the repository root enables workspace mode for seamless cross-module development:

\begin{verbatim}
go 1.21

use (
    ./star-proto/gen/go
    ./star-proto/tests/go
    ./star-gateway
)
\end{verbatim}
